% ============================================================
% PHẦN 1: TỔNG QUAN
% ============================================================
\section{Tổng quan}

% ------------------------------------------------------------
\subsection{Giới thiệu đề tài}

Trong thập kỷ gần đây, sự bùng nổ của Internet of Things (IoT) đã mang lại khả năng
giám sát liên tục và thu thập dữ liệu theo thời gian thực từ các hệ thống vật lý quy
mô lớn. Một trong những lĩnh vực ứng dụng có tầm quan trọng đặc biệt là quản lý mạng
cấp nước đô thị (Water Distribution Network --- WDN): hệ thống gồm hàng trăm cảm
biến đo mức nước, lưu lượng, áp suất và trạng thái bơm, liên tục phát sinh luồng dữ
liệu khổng lồ cần được thu thập, xử lý và phân tích kịp thời
\cite{taormina2018batadal}.

Tuy nhiên, việc xử lý và khai thác luồng dữ liệu quy mô lớn từ các hệ thống WDN đặt
ra nhiều thách thức kỹ thuật: khối lượng dữ liệu phát sinh liên tục (high-velocity),
yêu cầu độ trễ thấp (low-latency) trong phát hiện sự cố, và sự phức tạp trong việc
tích hợp nhiều thành phần xử lý phân tán. Các kiến trúc Big Data truyền thống như
Lambda Architecture thường phải duy trì song song cả Batch Layer lẫn Speed Layer,
dẫn đến sự trùng lặp trong logic xử lý. Kiến trúc Kappa \cite{kreps2014kappa} đề xuất
cách tiếp cận tinh gọn hơn: chỉ sử dụng một Speed Layer duy nhất (streaming-only),
giúp đơn giản hóa toàn bộ pipeline mà vẫn đáp ứng yêu cầu xử lý thời gian thực.

Bộ dữ liệu BATADAL (Battle of the Attack Detection ALgorithms) \cite{taormina2018batadal}
là một benchmark chuẩn được xây dựng nhằm đánh giá khả năng phát hiện bất thường
(anomaly detection) trên dữ liệu WDN. BATADAL cung cấp dữ liệu vận hành của hệ thống
C-TOWN --- một mạng nước đô thị mô phỏng gồm 43 cảm biến --- kèm theo nhãn tấn công
(attack label) tương ứng. Đây là nguồn dữ liệu lý tưởng để xây dựng và kiểm chứng
một hệ thống Big Data end-to-end: từ việc giả lập luồng dữ liệu IoT, xử lý streaming,
phát hiện sự cố, đến trực quan hóa kết quả theo thời gian thực.

Đề tài này xây dựng một pipeline phân tích dữ liệu lớn hoàn chỉnh dựa trên bộ dữ
liệu BATADAL, áp dụng kiến trúc Kappa với các công nghệ mã nguồn mở hiện đại:
Apache Kafka \cite{kreps2011kafka} làm message broker, Apache Spark Structured
Streaming \cite{zaharia2016spark} làm Speed Layer, InfluxDB làm time-series database,
và Grafana làm lớp trực quan hóa thời gian thực. Toàn bộ hệ thống được containerize
bằng Docker, hỗ trợ triển khai cả ở môi trường cục bộ lẫn trên nền tảng đám mây
(AWS EC2).

% ------------------------------------------------------------
\subsection{Ý nghĩa đề tài}

Đề tài mang lại ý nghĩa thiết thực trên nhiều phương diện:

\textbf{Về mặt kỹ thuật}, đề tài minh chứng cho khả năng xây dựng một hệ thống Big
Data hoàn chỉnh --- từ data ingestion, stream processing, storage đến visualization
--- hoàn toàn bằng các công nghệ mã nguồn mở, có thể triển khai trên một máy đơn
hoặc mở rộng lên cụm phân tán (distributed cluster). Kiến trúc được thiết kế theo
nguyên tắc containerization, đảm bảo tính tái lập (reproducibility) và dễ dàng mở
rộng quy mô (scalability).

\textbf{Về mặt ứng dụng}, hệ thống cung cấp khả năng giám sát mạng cấp nước theo
thời gian thực: phát hiện sớm các sự cố rò rỉ, tấn công mạng vật lý
(cyber-physical attack) hoặc hoạt động bơm bất thường thông qua cơ chế rule-based
anomaly detection dựa trên đặc tả kỹ thuật của mạng C-TOWN. Điều này có ý nghĩa
lớn trong bối cảnh cơ sở hạ tầng nước đô thị ngày càng được số hóa và trở thành
mục tiêu của các mối đe dọa an ninh mạng \cite{taormina2018batadal}.

\textbf{Về mặt học thuật}, đề tài vận dụng và kiểm chứng các khái niệm cốt lõi
trong lĩnh vực Big Data và Data Engineering: kiến trúc luồng dữ liệu (streaming
architecture), hợp đồng dữ liệu (data contract), phân tầng dữ liệu (Bronze/Silver
data tiers), kiểm thử hệ thống phân tán (unit test, integration test và end-to-end
test), và triển khai đám mây (cloud deployment) với các tiêu chuẩn bảo mật hiện đại
(container security scan bằng Trivy).

% ------------------------------------------------------------
\subsection{Phạm vi đề tài}

\subsubsection{Những nội dung thực hiện}

Trong khuôn khổ đề tài, nhóm thực hiện các nội dung sau:

\begin{enumerate}[label=\arabic*.]
  \item \textbf{Thiết kế kiến trúc hệ thống:} Xây dựng sơ đồ luồng dữ liệu theo
    kiến trúc Kappa với pipeline: IoT Simulator (Python) $\rightarrow$ Apache Kafka
    $\rightarrow$ Spark Structured Streaming $\rightarrow$ InfluxDB $\rightarrow$
    Grafana. Containerize toàn bộ hệ thống bằng Docker Compose, hỗ trợ cả môi
    trường local và cloud.

  \item \textbf{Định nghĩa tập dữ liệu cho bài toán:} Phân tích và lựa chọn tập dữ
    liệu phù hợp từ bộ BATADAL: sử dụng \texttt{BATADAL\_dataset03.csv} (8\,761 bản
    ghi, toàn bộ nhãn bình thường) làm nguồn lớp âm (negative class); xây dựng
    \texttt{BATADAL\_synthetic.csv} bổ sung các tình huống tấn công giả lập làm
    nguồn lớp dương (positive class). Loại bỏ \texttt{BATADAL\_dataset04.csv} khỏi
    quá trình huấn luyện do nhãn \texttt{ATT\_FLAG\,=\,-999} (dữ liệu mù, chưa gán
    nhãn). Lựa chọn 31 đặc trưng vật lý (mức nước bể $L_{Ti}$, lưu lượng bơm
    $F_{PUj}$, trạng thái đóng/mở $S_{PUj}$, $S_{V2}$), loại bỏ đặc trưng thời
    gian do thuật toán time compression làm mất ý nghĩa tuyệt đối của timestamp.

  \item \textbf{Xây dựng IoT Data Simulator:} Viết Python producer đọc dữ liệu từ
    file CSV (BATADAL dataset03/04), áp dụng thuật toán time compression để giả lập
    luồng dữ liệu real-time với tốc độ tối thiểu 50 messages/giây, gửi vào Kafka
    topic \texttt{water\_stream}.

  \item \textbf{Xây dựng Speed Layer với Spark Structured Streaming:} Consume dữ
    liệu từ Kafka, thực hiện tiền xử lý (lọc null, ép kiểu), áp dụng rule-based
    anomaly detection dựa trên thông số vận hành của mạng C-TOWN, ghi kết quả vào
    InfluxDB (measurement \texttt{water\_telemetry}).

  \item \textbf{Lưu trữ và trực quan hóa:} Sử dụng InfluxDB làm time-series database
    để lưu trữ dữ liệu telemetry và cờ cảnh báo. Xây dựng Grafana dashboard hiển thị
    trạng thái mạng nước theo thời gian thực, bao gồm mức nước bể chứa, trạng thái
    bơm, áp suất và cảnh báo sự cố.

  \item \textbf{Huấn luyện mô hình ML phát hiện bất thường (offline):} Xây dựng
    module \texttt{train\_model.py} huấn luyện mô hình Random Forest Classifier
    (100 cây, \texttt{class\_weight=``balanced''}) trên tập dữ liệu kết hợp
    dataset03 và synthetic. Đánh giá mô hình theo chiến lược stratified 75/25 split,
    tính Precision, Recall, F1-score và in confusion matrix. Mô hình huấn luyện
    trên toàn bộ dữ liệu được lưu ra \texttt{anomaly\_model.pkl}. Tiếp theo,
    module \texttt{stream\_processor.py} tải mô hình, phân phối đến các executor
    thông qua Spark broadcast variable, và gọi inference trên từng micro-batch
    bằng \texttt{pandas\_udf}, ghi kết quả dự đoán vào trường \texttt{ml\_alert}
    trong InfluxDB song song với cơ chế rule-based, cho phép so sánh hai nguồn
    cảnh báo trên cùng một Grafana dashboard.

  \item \textbf{Kiểm thử hệ thống:} Viết và thực thi đầy đủ unit tests, integration
    tests và end-to-end (E2E) tests với code coverage $\geq 100\%$. Thực hiện container
    security scan bằng Trivy nhằm đảm bảo không có lỗ hổng bảo mật mức CRITICAL.

  \item \textbf{Triển khai đám mây:} Triển khai các dịch vụ backend (Kafka, InfluxDB,
    Grafana) lên AWS EC2 (Singapore region, Ubuntu 24.04), cấu hình Security Group và
    tường lửa UFW, cho phép Spark driver và producer kết nối từ máy local vào cloud
    theo mô hình hybrid deployment.
\end{enumerate}

\subsubsection{Những nội dung không thực hiện}

Trong phạm vi của đề tài, nhóm không thực hiện các nội dung sau:

\begin{enumerate}[label=\arabic*.]
  \item \textbf{Batch Layer:} Đề tài theo kiến trúc Kappa (streaming-only), do đó
    không xây dựng Batch Layer để tái xử lý dữ liệu lịch sử theo mô hình Lambda
    Architecture.

  \item \textbf{Infrastructure as Code (IaC):} Đề tài không sử dụng các công cụ
    như Terraform hay Ansible để tự động hóa việc khởi tạo và cấu hình hạ tầng
    đám mây.

  \item \textbf{Hệ thống cảnh báo ngoài Grafana:} Đề tài không tích hợp các kênh
    thông báo bên ngoài (email, Slack, PagerDuty, v.v.) khi phát hiện sự cố. Cảnh
    báo chỉ được hiển thị trực quan trong Grafana dashboard.

  \item \textbf{Bộ dữ liệu điện (Electricity dataset):} Mặc dù bộ dữ liệu
    \textit{Electricity Load Diagrams 2011--2014} (UCI) có trong thư mục dataset,
    đề tài này chỉ tập trung vào bộ dữ liệu nước BATADAL. Bộ dữ liệu điện không
    được sử dụng trong pipeline.
\end{enumerate}

% ------------------------------------------------------------
\subsection{Bố cục bài báo cáo}

Bài báo cáo được tổ chức thành bảy phần như sau:

\textbf{Phần 1 --- Tổng quan} (phần hiện tại): Giới thiệu bối cảnh đề tài, ý nghĩa,
phạm vi thực hiện và bố cục báo cáo.

\textbf{Phần 2 --- Các công trình nghiên cứu liên quan}: Trình bày tổng quan các
nghiên cứu về phát hiện bất thường trong mạng cấp nước, kiến trúc Big Data cho
dữ liệu IoT, và các hệ thống giám sát thời gian thực.

\textbf{Phần 3 --- Cơ sở lý thuyết}: Trình bày các kiến thức nền tảng về kiến trúc
Kappa, Apache Kafka, Spark Structured Streaming, InfluxDB, Grafana, và bộ dữ liệu
BATADAL.

\textbf{Phần 4 --- Phân tích vấn đề}: Phân tích yêu cầu hệ thống, đặc điểm dữ liệu
BATADAL, và các thách thức kỹ thuật khi xây dựng pipeline phân tích dữ liệu lớn
cho IoT mạng nước.

\textbf{Phần 5 --- Giải pháp đề xuất và hiện thực}: Mô tả chi tiết kiến trúc hệ
thống, thiết kế từng component trong pipeline, và cách triển khai trên môi trường
Docker và AWS EC2.

\textbf{Phần 6 --- Thực nghiệm}: Trình bày kết quả kiểm thử (unit, integration,
E2E), kết quả phát hiện bất thường trên bộ dữ liệu BATADAL, và đánh giá hiệu
năng hệ thống (throughput, latency).

\textbf{Phần 7 --- Kết luận và hướng phát triển}: Tổng kết kết quả đạt được, nhận
xét về các hạn chế và đề xuất các hướng mở rộng trong tương lai.
